\section{Verificación de implementación y mejora continua}
\subsection{Lista de verificación habitual}
Graham recomienda ejecutar una lista de verificación antes y después de cada iteración del Agent, para asegurar que el prompt, la cadena de herramientas, los registros y el mecanismo de aprobación estén en condiciones de uso. Los puntos de la lista de verificación incluyen: si se actualizaron los open-source requirements, si se verificaron los permisos de la imagen, si se actualizó la línea base de pruebas, si se sincronizó la notificación de Slack.

\begin{table}[H]
\centering
\begin{tabular}{p{0.32\textwidth} p{0.6\textwidth}}
\toprule
Punto de verificación & Descripción \\
\midrule
Estado del prompt & Asegurar que la plantilla de prompt incluya las muestras de issue más recientes y los guardrails \\
\midrule
Actualización de la cadena de herramientas & Verificar si la imagen de docker, las versiones de dependencias y los permisos de la API están vencidos \\
\midrule
Registros y auditoría & Verificar que el flujo de eventos sea normal y que el audit log pueda reproducirse \\
\midrule
Mecanismo de aprobación & Revisar si el proceso de aprobación de comandos de alto riesgo está completo, incluyendo el emergency override \\
\bottomrule
\end{tabular}
\caption{Puntos de verificación habituales antes y después de la implementación.}
\end{table}

\subsection{Ciclo de mejora continua}
Cada iteración del Agent debe incluir una etapa de revisión: revisar los casos exitosos y las experiencias de fracaso, y escribir los nuevos insight en el FAQ o la documentación, para usarlos en la siguiente ronda de entrenamiento del modelo. Graham enfatiza la «normalización de la mejora del 1%»: cada día mejorar un poco alguno de los elementos entre las instrucciones, los prompts y la cadena de herramientas.

\begin{importantbox}{Pasos pequeños, avance rápido}
Descomponer las actualizaciones del Agent en cambios de 5 minutos, y fusionarlos solo después de que pasen la verificación, evitando en la medida de lo posible los cambios masivos de una sola vez.
\end{importantbox}

\begin{knowledgebox}{Indicadores de mejora}
Usar la «proporción de casos exitosos», la «tasa de finalización de las retrospectivas» y el «número de actualizaciones de documentación» como indicadores de mejora, para cerrar el ciclo entre el trabajo de ingeniería y el de investigación.
\end{knowledgebox}

\subsection{Resumen de la sección}
\begin{itemize}
    \item La lista de verificación garantiza que la arquitectura del Agent esté en buen estado antes y después de cada iteración.
    \item Los pasos pequeños y rápidos junto con los indicadores de mejora permiten que el equipo progrese de manera continua bajo un riesgo controlado.
\end{itemize}

